\begin{otherlanguage}{italian}
\chapter*{Sommario}

\begin{spacing}{1.0}
Il rilevamento delle intrusioni di rete si affida ai sistemi a firme, che non vedono gli attacchi
senza una firma nota, e ai rilevatori basati sui dati, fragili quando il traffico cambia. Gli
approcci neurosimbolici promettono di unire le due strade portando la conoscenza di dominio
nell'apprendimento. Questa tesi costruisce un benchmark che confronta le Logic Tensor Networks con
baseline simboliche, di machine learning classico e di deep learning su NSL-KDD, CIC-IDS2017 e
CIC-DDoS2019, con la stessa rete delle controparti solo neurali e lo stesso budget di tuning. Misura
qualità del rilevamento, costo e usabilità, anche sotto stress: poche etichette, classi rare,
attacchi mai visti, trasferimento tra dataset, conoscenza degradata, rumore ed evasione. Due
metriche a livello di regola la completano.

Su questi dataset e basi di conoscenza la maggior parte dei risultati è negativa. La conoscenza
usata come vincolo in addestramento non dà un guadagno dimostrato: il suo effetto sul macro-F1 è di
$+0{,}008$, $+0{,}015$ e $+0{,}002$, non significativo su nessun dataset dopo la correzione di Holm.
Neanche con poche etichette aiuta: dove l'effetto si distingue dai seed è negativo, al massimo 1,5
punti. Le regole estratte dai dati non reggono fuori dall'addestramento: sui due dataset con un test
separato la precisione scende a 0,130 e 0,547, mentre le regole di meccanismo delle stesse basi
restano tra 0,950 e 0,984. Queste contano in inferenza in un solo caso, il trasferimento da
CIC-IDS2017 a CIC-DDoS2019: i modelli appresi non superano un macro-F1 di 0,4812, mentre la regola
sul flood UDP del dataset di destinazione, da sola, arriva a 0,918; in addestramento la stessa
regola porta la rete solo al livello della migliore baseline neurale. In ogni dataset, far correggere
le predizioni dall'intera base di conoscenza peggiora ogni modello appreso. In inferenza una rete
addestrata con la logica costa come una rete normale, e ogni assioma aggiunge dal 2 al 4\,\% al
tempo di un'epoca di addestramento. Su CIC-IDS2017 gli ensemble di alberi battono ogni rete. Il
contributo pratico della tesi è una checklist che dice quando conviene usare un approccio
neurosimbolico e quando no.
\end{spacing}
\end{otherlanguage}
